% Síntesis de la distinción institucional y contable de FPC, fcp_gg.tex.
\begin{frame}[t]{PGN, GNC y gobierno general no son lo mismo}
\small
\noindent\begin{tabularx}{\textwidth}{@{}>{\raggedright\arraybackslash}p{3.4cm}>{\raggedright\arraybackslash}X>{\raggedright\arraybackslash}p{3.3cm}@{}}
\toprule
\textbf{Concepto} & \textbf{Cobertura} & \textbf{Lectura} \\
\midrule
\textbf{PGN}\par
Presupuesto General de la Nación & Presupuesto nacional y presupuestos de los establecimientos públicos nacionales. & Autoriza gastos de funcionamiento, inversión y deuda. \\
\addlinespace[0.55em]
\textbf{GNC}\par
Gobierno Nacional Central & Operaciones fiscales del nivel central. No equivale a todo el PGN. & Balance fiscal y regla fiscal nacional. \\
\addlinespace[0.55em]
\textbf{Gobierno general} & Gobierno central, gobiernos regionales y locales, y seguridad social. Se consolidan las transferencias internas. & Balance consolidado de esos subsectores. \\
\bottomrule
\end{tabularx}\par
\vspace{0.35cm}
\noindent\textbf{PGN y GNC: conciliar cobertura, clasificación y momento de registro.}
\note{[Sources] FPC, modules/fcp_gg.tex: esquema de desagregación, balance del SPNF y distinción entre contabilidad presupuestal y fiscal. Decreto 111 de 1996, artículos 3 y 11: https://www1.funcionpublica.gov.co/eva/gestornormativo/norma.php?i=5306.
MHCP, Balance del Gobierno General: https://www.minhacienda.gov.co/politica-fiscal/cifras-pol%C3%ADtica-fiscal/gobierno-general/balance.
MHCP--CIEFP, Clasificación de entidades del sector público colombiano, versión 4, salvedad sobre cobertura por transacciones en estadísticas tradicionales: https://www.minhacienda.gov.co/documents/20119/1528668/4.%2BClasificaci%C3%B3n%2Bde%2Bentidades%2Bdel%2Bsector%2Bp%C3%BAblico%2Bcolombia_V4.pdf/7d1512fb-7139-12db-bede-33f7112a52e3?download=true.
Anexo al Mensaje PGN 2027 reformulado, agosto de 2026, sección 2.4, páginas impresas 204--206, cuadros 2.4.2--2.4.6. El gobierno central comprende GNC y otros componentes; no se identifica PGN con GNC ni con todo el gobierno general. Consulta: 2 de septiembre de 2026.}
\end{frame}
